\section{Introduction}
\label{sec:intro}

\begin{figure}[t]
\centering
\resizebox{0.97\linewidth}{!}{%
\begin{tikzpicture}[
  font=\small,
  box/.style={draw, rounded corners=2pt, align=center, inner sep=3pt, minimum height=0.8cm},
  arr/.style={-{Stealth[length=2mm]}, thick}
]
% ---------------- timeline ----------------
\node[anchor=west, font=\small\bfseries] at (0,2.95) {(a) Labels from a single \qwen trajectory};
\fill[green!18] (0,1.5) rectangle (3.0,2.1);
\fill[gray!12] (3.0,1.5) rectangle (7.2,2.1);
\fill[red!25] (7.2,1.5) rectangle (12.0,2.1);
\draw (0,1.5) rectangle (12.0,2.1);
\node at (1.5,1.8) {safe prefix $x_{\le t}$};
\node at (5.1,1.8) {harmless background};
\node at (9.6,1.8) {actionable harm};
\draw[very thick, blue!70!black] (3.0,1.35) -- (3.0,2.25);
\node[blue!70!black, above] at (3.0,2.25) {cutoff $t$};
\draw[very thick, red!70!black] (7.2,1.35) -- (7.2,2.25);
\node[red!70!black, above] at (7.2,2.25) {onset $\onset$};
% horizon brackets
\foreach \h/\x/\lab/\y in {64/4.0/{$y_{64}{=}0$}/1.0, 128/5.0/{$y_{128}{=}0$}/0.65, 256/7.0/{$y_{256}{=}0$}/0.3, 512/11.0/{$y_{512}{=}1$}/-0.05}{
  \draw[|-|] (3.0,\y) -- (\x,\y);
  \node[anchor=west, font=\scriptsize] at (\x+0.05,\y) {\lab};
}
\node[anchor=west, font=\scriptsize, align=left] at (12.2,1.8) {$y_H{=}1$ iff\\$\onset < t{+}H$};
% ---------------- observers ----------------
\node[anchor=west, font=\small\bfseries] at (0,-0.55) {(b) Observers read frozen activations of $x_{\le t}$ only};
\node[box, fill=gray!8] (acts) at (1.4,-2.35) {residuals\\L9 / L18 / L27};
\node[box, fill=yellow!15] (probe) at (5.0,-1.6) {linear probe\\(+16-token history)};
\node[box, fill=orange!15] (ao) at (5.0,-2.6) {off-the-shelf AO\\``harmful later?''};
\node[box, fill=blue!10] (direct) at (5.0,-3.65) {\direct reader\\encoder $\to$ memory slots};
\node[box, fill=purple!12] (fc) at (9.3,-3.15) {\forecast reader\\4 horizon queries $\to$\\forecast vectors $\to$ timing head};
\node[box, fill=white] (qa) at (9.3,-4.35) {shared QA head\\(question embedding)};
\node[box, fill=green!10] (out) at (13.0,-2.35) {$\hat p(\onset<t{+}H)$\\early-warning alert};
\draw[arr] (acts) -- (probe);
\draw[arr] (acts) -- (ao);
\draw[arr] (acts.south) |- (direct.west);
\draw[arr] (direct.east) -- ++(0.5,0) |- (qa.west);
\draw[arr] (direct.east) -- ++(0.5,0) |- (fc.west);
\draw[arr] (fc) -- (qa);
\draw[arr] (probe.east) -- (out.west);
\draw[arr] (ao.east) -- (out.west);
\draw[arr] (qa.east) -| (out.south);
\end{tikzpicture}}
\caption{{\bf Task and observers.} (a) From one \qwen generation we take still-safe prefixes at fixed cutoffs $t\in\{32,\dots,256\}$. The binary label $y_H$ says whether the first clear violation (onset $\onset$, located by an LLM annotator) falls within the next $H$ tokens. (b) All observers see only activations of the truncated text. The \direct and \forecast readers share the same encoder, QA head, labels, and training budget. The \forecast reader must pass all information through horizon-specific forecast vectors and a censoring-aware timing head.}
\label{fig:overview}
\end{figure}


Most safety monitors for language models answer one question: is the text produced so far unsafe? This question comes too late for an important class of failures. A model asked for a ``long-form overview'' of a dangerous topic can spend hundreds of tokens on history and general context, and only then move into step-by-step instructions. A monitor that judges only the current text raises an alarm after the harmful content is already on the screen. A monitor that can \emph{forecast} the turn could stop, steer, or escalate the generation before any harm is emitted.

There is good reason to think such forecasts are possible. Hidden states carry information about future tokens~\citep{pal2023futurelens}, about global properties of the coming response~\citep{dong2025emergent}, and even about how many tokens remain~\citep{merzouk2026remaining}. Activation verbalizers make this information readable. LatentQA~\citep{pan2024latentqa} and Activation Oracles (AOs)~\citep{karvonen2025activation} fine-tune a copy of the model to answer natural-language questions about another model's activations. This makes them a natural candidate for asking a model ``will this response turn harmful later?''

However, existing work does not test this. Verbalizers are evaluated on questions about the \emph{current} state: hidden personas, secret words, or the next few tokens~\citep{karvonen2025activation,bauer2026building}. Streaming guards mostly detect harm already present in the prefix~\citep{li2025judgment,sharma2025constitutional}, and the one forecasting guard we know of is text-only and looks only a short distance ahead~\citep{kavumba2026predict}. On the other hand, recent studies warn that trained AOs at best tie linear probes~\citep{bauer2026building}. Any claim that ``verbalizers forecast'' must therefore beat probes and text baselines on identical inputs. Finally, if a reader trained on future labels does well, it is unclear whether the gain comes from the \emph{supervision} or from an \emph{architecture} that explicitly represents the future.

We study both questions on natural generations of \qwen (\figref{fig:overview}). {\bf Can the activations of a still-safe prefix tell us that a clear safety violation will first appear within the next $H$ tokens?} We build a family-split benchmark of 3,520 safe prefixes from 880 prompt families, with token-level onsets from an LLM annotator and labels that account for censoring. We compare prevalence, position, and text baselines, activation probes, \qwen self-report, and an off-the-shelf AO. {\bf Does an explicit future representation help?} We train two readers with the same activations, labels, QA head, and budget. One is a direct AO-style reader. The other routes everything through horizon queries, forecast vectors, and a censoring-aware timing head.

Delayed onsets are partly forecastable. A linear probe on the last token reaches \auroc 0.74 at $H{=}64$ and 0.87 at $H{=}512$. Adding the mean of the last 16 layer-18 activations raises this to 0.79 and 0.90. At a 5\% benign false-alarm rate, this history probe flags 56\% of unsafe trajectories at least 128 tokens before onset, against 42\% for a single-position probe or a prompt+prefix text model. The off-the-shelf AO is the weakest non-trivial method: it says ``No'' to about 99\% of prefixes and scores \auroc 0.54 at $H{=}64$. The forecast representation does not beat the matched direct reader. Its timing head trails the direct reader by 0.02--0.04 \auroc, with family-paired bootstrap intervals that exclude zero at $H{\ge}256$. Both trained readers are no better than a linear probe with history.

Our contributions are:
\begin{itemize}[leftmargin=*,itemsep=0pt,topsep=0pt]
    \item We construct a lead-distance benchmark for delayed unsafe behavior on fresh \qwen generations. It is split by prompt family before generation, uses token-level onsets, and checks that observer inputs are causal.
    \item We measure how forecastable onsets are at 64--512 tokens for text baselines, probes, self-report, an off-the-shelf AO, and trained readers. Activation history is the single most useful addition, and an off-the-shelf AO is strongly biased toward ``safe''.
    \item We run a matched comparison between a direct AO-style reader and a forecast-representation reader. The explicit forecast adds nothing, so the gain over zero-shot verbalizers comes from future-event supervision.
    \item We complement the main results with shortcut controls (shuffled and constant activations, prompt-only text), a readout-ceiling diagnostic, a held-out question meaning, and a separate fictional-harm setting.
\end{itemize}

\secref{sec:related} reviews related work. \secref{sec:method} describes the benchmark and the observers. \secref{sec:results} reports results, and \secref{sec:discussion} discusses their meaning and limits.
